\bookchapter{第四十八章：認識安卓——從操作系統到開發生態}

\chapterintro{恭喜汝做出此決定——從「玩手機之人」變成「做手機應用之人」。於動手之前，吾儕先將地圖攤開：安卓從哪裡來？為何全世界之開發者皆於為它寫應用？一個App從汝之手指敲鍵盤到用戶手機上，中間皆發生了何？}

\booksection{一、安卓之身世：從「小機器人」到「全球最大系統」}

安卓（Android）之故事始於2003年——一個名叫Andy Rubin之工程師創辦了Android公司，目標甚樸素：做一套「更懂相機」之智能操作系統。後來谷歌於2005年收購了它，2007年聯合幾十家廠商成立開放手機聯盟（OHA），2008年9月，第一款安卓手機HTC Dream（G1）正式發售。

後面之故事大家皆知道了：到2026年，安卓乃全球市場份額最大之行動操作系統，活躍設備數以十億計。汝之手機、電視盒子、車載大屏、智能手錶，甚可能皆於跑安卓。此意味著何？意味著汝所寫之每一行代碼，皆有機會被世界上最多之人用到。

\begin{tcbinfo}{小訣：安卓之開源基因}
安卓之核心代碼（AOSP，Android Open Source Project）乃開源之——任何廠商皆能拿去修改定製。故小米有MIUI（而今叫澎湃OS）、華為有鴻蒙兼容層、三星有One UI。此對開發者乃好事：汝學之乃「安卓開發」，不為「某個品牌之開發」，一套功夫，全網通喫。

\end{tcbinfo}

版本號之進化也甚有意思：從1.5之紙杯蛋糕（Cupcake）開始，每個大版本皆以甜點命名——甜甜圈、閃電泡芙、凍酸奶、薑餅……一直到2019年之安卓10才改成了純數字。到2026年8月，最新正式版乃Android 17，而開發時汝真正須關心之「最低支持版本」，主流App一般定於Android 8（API 26）到Android 10（API 29）之間——支持到幾年前之設備，汝之用戶盤子才夠大。

\begin{center}
\begin{tabularx}{\linewidth}{|X|X|X|}
\hline
\textbf{版本} & \textbf{代號／年份} & \textbf{對開發者重要之事} \tabularnewline
\hline
Android 8 (API 26) & Oreo · 2017 & 通知渠道、後臺限制——甚多老機型之最低線 \tabularnewline
Android 10 (API 29) & 2019 & 分區存儲登場，文件訪問規矩大改 \tabularnewline
Android 12 (API 31) & 2021 & Material You動態取色，界面觀感大升級 \tabularnewline
Android 13 (API 33) & 2022 & 運行時通知權限，彈窗問汝「準不準發通知」 \tabularnewline
Android 14-15 (API 34-35) & 2023-2024 & 預測性返回、16KB頁、隱私進一步收緊 \tabularnewline
Android 16-17 (API 36-37) & 2025-2026 & 更快之小版本節奏、AI能力整合（如Gemini擴展） \tabularnewline
\hline
\end{tabularx}
\end{center}

\begin{center}\small\color{gray}安卓版本簡史（截至2026年8月，開發重點關注項）\end{center}

\booksection{二、用何語言寫：Java抑Kotlin？}

安卓開發歷史上出現過兩種主流語言。Java乃「開國元老」——安卓誕生之初就用它；Kotlin乃2017年穀歌欽點之「一等公民」，2019年起官方明確「Kotlin優先」（Kotlin-first），到2026年，新項目幾乎全部用Kotlin，舊項目也大多於陸續遷移。

說句實於話：若汝乃完全零基礎，直接學Kotlin，勿回頭碰Java。原因有三：一為Kotlin代碼量比Java少30\%-40\%（同樣之功能，Kotlin更短更清晰）；二為它天生防「空指針」此種新手噩夢；三為所有現代教程、官方文檔、AI編程工具（還記得吾儕之AI編程書嗎？）皆於以Kotlin為準。Java只於維護老項目、讀老代碼時會遇到。

Java風格（囉嗦版）

\begin{lstlisting}
public class MainActivity extends Activity \{
@Override
protected void onCreate(Bundle b) \{
super.onCreate(b);
setContentView(R.layout.activity);
\}
\}
\end{lstlisting}

Kotlin風格（清爽版）

\begin{lstlisting}
class MainActivity : ComponentActivity() \{
override fun onCreate(b: Bundle?) \{
super.onCreate(b)
setContent \{ App() \}
\}
\}
\end{lstlisting}

同樣一件事，Kotlin少了一堆花括號與樣板。第50章吾儕會正式學它，而今汝只須記住：安卓開發之「普通話」叫Kotlin。

\booksection{三、一個App為何以誕生之？}

甚多人以為「開發App」即打開電腦狂寫代碼。實際上，一個App從想法到用戶手機上，大致要經過如此一條流水線：

想法

設計（界面+功能）

寫代碼

編譯打包（APK/AAB）

測試（模擬器+真機）

發佈（商店/安裝包）

\begin{center}\small\color{gray}圖1：一個安卓App從想法到上架之全流程\end{center}

流程裡有兩個概念汝而今就要認識：

\begin{itemize}
  \item APK：安卓應用之安裝包格式（.apk），雙擊就能裝到手機上，就像Windows之.exe。平時自己測試、予朋友傳著玩，用它。
  \item AAB：Android App Bundle（.aab），上架Google Play用之「高級打包格式」。商店會根據用戶之手機型號動態裁剪，使每個人下載之包更小。自己裝不了，只能交予商店。
\end{itemize}

另外汝還會不斷聽到一個詞：SDK（Software Development Kit，軟件開發工具包）。它即安卓官方提供之一整套「開發裝備」——包括各版本之API庫、調試工具、模擬器鏡像等。汝寫代碼時調用之Activity、Button此些類，皆來自SDK。

\booksection{四、安卓技術棧全景：汝之「武器清單」}

正式學習前，先認識一下2026年主流安卓項目之標準配置——此張清單汝而今看不懂沒關係，全書學完回來看，會豁然開朗：

\begin{center}
\begin{tabularx}{\linewidth}{|X|X|X|}
\hline
\textbf{層次} & \textbf{主流選擇（2026）} & \textbf{所用為何} \tabularnewline
\hline
語言 & Kotlin 2.x & 寫邏輯之語言 \tabularnewline
界面 & Jetpack Compose + Material 3 & 用代碼畫界面（本書主線） \tabularnewline
架構 & MVVM + ViewModel & 將界面與邏輯分開，代碼不亂 \tabularnewline
數據存儲 & Room（數據庫）＋DataStore（配置） & 將數據存進手機 \tabularnewline
網絡 & Retrofit＋OkHttp＋Kotlin協程 & 與服務器對話 \tabularnewline
構建工具 & Gradle + Kotlin DSL & 將代碼變成安裝包 \tabularnewline
開發環境 & Android Studio & 谷歌官方之「全家桶」IDE \tabularnewline
\hline
\end{tabularx}
\end{center}

\begin{center}\small\color{gray}2026年安卓主流技術棧速覽\end{center}

勿被此一串英文嚇到。它等就像廚房裡之鍋碗瓢盆——汝而今只須知道「做飯要進廚房」，具體每樣工具何以用，後面章節逐個拆予汝看。而且好消息為：Android Studio創建新項目時會自動幫汝將此一整套配置好，汝只管炒菜。

\booksection{五、安卓vs iOS：開發者之第一堂「選邊課」}

汝手裡可能正同時有安卓與蘋果設備。作為開發者，兩邊之體驗差異值得提前瞭解——不為「哪個更好」之站隊，乃「哪邊更適合汝起步」之現實賬：

\begin{center}
\begin{tabularx}{\linewidth}{|X|X|X|}
\hline
\textbf{維度} & \textbf{安卓開發} & \textbf{iOS開發} \tabularnewline
\hline
開發電腦 & Windows／Mac／Linux皆行 & 必須Mac（Xcode獨佔） \tabularnewline
開發工具 & Android Studio（免費） & Xcode（免費）＋上架年費99美元 \tabularnewline
上架審核 & 寬鬆，幾小時至幾天 & 嚴格，1-7天乃常態 \tabularnewline
裝機試玩 & APK隨便發、隨便裝 & 分發受限，測試須簽名／TestFlight \tabularnewline
收入環境 & 廣告／國內渠道為主 & App Store內購付費習慣更好 \tabularnewline
學習成本 & 概念較多（四大組件等） & 生態封閉但體系統一、文檔精良 \tabularnewline
\hline
\end{tabularx}
\end{center}

\begin{center}\small\color{gray}安卓與iOS開發路徑對比（2026年8月）\end{center}

說句實於話：若汝之電腦不為Mac，那此題根本不用選——安卓乃唯一選項。若汝有Mac，本書第八、九兩篇可對照著讀，編程思路為相通之，學完一邊再學另一邊會快甚多。

\begin{tcbpractice}{動手一試：看看汝手機裡之「安卓」}
打開手機之［設定］［關於手機］（或「關於本機」），尋到「安卓版本」或「Android版本」那一欄——連續點它幾下，還能觸發系統之彩蛋小遊戲。記下汝之版本號，對照上面之版本表看看它屬於哪一代。汝之手機，即汝之第一個「開發測試機」。

\end{tcbpractice}

\begin{tcbsummary}{本章小結}
\begin{itemize}
  \item 安卓2008年誕生，2026年乃全球市場份額最大之行動系統，活躍設備數十億
  \item 最新正式版Android 17；主流App最低支持Android 8-10，用戶盤子才夠大
  \item 新項目一律用Kotlin：代碼少30\%以上、防空指針、官方欽定
  \item APK自己裝，AAB上商店；SDK乃官方開發裝備包
  \item 2026標準技術棧：Kotlin＋Compose＋MVVM＋Room＋Retrofit＋Gradle＋Android Studio
\end{itemize}

\end{tcbsummary}
